% !TeX root = ../../Book3.tex
% 纲要标题：## 附录 C　奇点与交换代数专题
% 纲要位置：工作记录/计划纲要.md:2908

\chapter{奇点与交换代数专题}
\label{b3:app:C}
\BookChapterNumber{b3:app:C}

\begin{chapterabstract}
本附录通过几个可计算模型，介绍交换代数在奇点研究中的进一步用途。形式纤维说明完备化需要比较哪些点；Frobenius 的自由性与分裂使正特征运算转为模论问题；完全交的余切与切复形保留方程及其变化；局部交数、乘数理想和多次数对偶则分别提供相交、消失阶与同调对偶方面的信息。
\end{chapterabstract}

\begin{learninggoals}
\begin{itemize}
\item 计算简单局部环的闭形式纤维与泛形式纤维，说明它们检验了卓越性的哪一项条件。
\item 分清 Frobenius 的有限、平坦与分裂性质，并检验具体算例。
\item 在完全交模型中区别余切复形与切复形，计算超曲面的一阶方程扰动。
\item 用商环长度计算局部交数，在正规交叉模型中计算乘数理想，并比较不同次数的 Ext。
\end{itemize}
\end{learninggoals}

交叉直线 \(xy=0\) 与尖点 \(y^2=x^3\) 在原点的局部环都为一维，嵌入维数却为二：它们都有单元素参数系，但极大理想不能由一个元素生成，因而不是正则局部环。二者都有二重的初始方程，前者却有两条分支，后者只有一条。单用“这是奇点”或单用一个重数，无法解释这些差别。要继续研究，须问完备化在各素点上增加了什么、正特征下 Frobenius 怎样作用，以及方程之间是否还有隐藏的关系。

\ProseParagraphBreak

这四节从不同方面研究奇点，并不构成依次推出彼此的一条定理链。形式纤维比较原环与完备化的点，余切与切复形保留方程及其关系，局部交数衡量分支相遇的程度；正特征和特征零又分别提供 Frobenius 与乘数理想的方法。本附录用可计算模型逐一说明工具解决的问题，并在需要深理论的地方明确指出所用结论的范围。

\ProseParagraphBreak

所需基础为本卷的局部化、完备化、正则序列与自由消解。\secref{b3:sec:C03}还用到第14章的微分模及第23章的平方零提升；\secref{b3:sec:C04}的末段调用第22章已经证明的局部对偶。卓越环的统一定义在附录 D 的\cref{b3:def:excellent-ring}，本附录先回顾其三个条件。余切复形采用显式完全交模型；乘数理想还使用对数消解与相对典范除子，其存在性及选择独立性采用文献定理，计算则落在正规交叉与两个吹起坐标图中。

% 本附录按既有四节组织；具体模型给出证明，长证明给出构造概要及出处。

\cref{b3:tab:C-tools-and-hypotheses}比较这些工具实际保留的信息与本附录所用的假设。表中的余切与切复形列出完全交模型的范围，局部对偶列出给定正则局部覆盖的范围。

\begin{table}[htbp]
\centering
\caption{奇点研究工具、所测信息与适用条件}
\label{b3:tab:C-tools-and-hypotheses}
\begin{tabularx}{\linewidth}{@{}>{\raggedright\arraybackslash}p{0.15\linewidth}>{\raggedright\arraybackslash}p{0.27\linewidth}>{\raggedright\arraybackslash}p{0.29\linewidth}>{\raggedright\arraybackslash}X@{}}
\toprule
工具 & 所测信息 & 适用条件 & 本附录的证明范围 \\
\midrule
形式纤维 & 完备化在各素点上增加的信息 & Noether 局部环，几何正则性须检查域扩张 & 简单纤维完整计算；卓越性回指附录 D \\
Frobenius & 扭曲标量作用后的有限、平坦与分裂性质 & 特征 \(p>0\)；多项式自由基使用完美域 & 自由基与分裂算例证明；Kunz 判据引用 \\
余切复形 & 微分、方程及其高阶关系 & 特征零域上的正则序列商 & 完全交模型引用，DG 结构及同调展开 \\
切复形 & 导子与坐标变化不能消去的扰动 & 完全交导出对偶；固定超曲面展示 & 方程扰动完整计算；抽象分类引用 \\
局部交数 & 两条平面曲线相交的厚度 & \(k[\![x,y]\!]/(f,g)\) 有限长度 & 有限长度及 Tor 模型完整计算 \\
乘数理想 & 除子消失阶与相对典范除子的修正 & 光滑复簇、非零凝聚理想、正有理参数 & 定义及坐标算例；存在性与独立性引用 \\
对偶复形 & 非 CM 商环在多个次数的对偶信息 & 完备正则局部环及其非零商环 & 第22章结果调用，显式对偶模型计算 \\
\bottomrule
\end{tabularx}
\end{table}

% !TeX root = ../../Book3.tex
% 纲要标题：### C.1 卓越性、形式纤维与普适链性质
% 纲要位置：工作记录/计划纲要.md:2867

\section{卓越性、形式纤维与普适链性质}
\label{b3:sec:C01}

% 纲要内容（仅作写作计划，尚非教材正文）：
% * 形式纤维及闭点、泛点模型；G环的平坦完成映射与几何正则纤维，普适链性质及J-2的分工
% <!-- 短节：不设小节 -->


% 正文以具体模型说明专题；长证明给出概要并指明技术细节的出处。

对一个局部环取完备化，常能把多项式计算改成幂级数计算。但忠实平坦本身并没有说明每条纤维的几何性质，也没有保证所有有限型代数的正则点都组成开集。研究奇点时，往往还希望完成下列工作：在完备化前后比较局部结构；在一个族中找出统一的正则开集；沿不同素理想链计算一致的余维。卓越性把这些需求组织为明确条件。

\ProseParagraphBreak
设 \((A,\mathfrak m)\) 为 Noether 局部环，\(\widehat A\) 为其 \(\mathfrak m\) 进完备化。对 \(\mathfrak p\in\operatorname{Spec}A\)，环
\[
\widehat A\otimes_A\kappa(\mathfrak p)
\]
称为该素点处的\term{形式纤维环}[formal fibre ring]，其中 \(\kappa(\mathfrak p)=\operatorname{Frac}(A/\mathfrak p)\) 是该素点的剩余域。闭纤维为 \(A/\mathfrak m\)，其他素点的纤维则可能包含新的信息。相关定义见松村英之《Commutative Ring Theory》\cite[\S 32, pp. 255--256]{Matsumura1989CommutativeRingTheory}。

形式纤维通常不是剩余域上的有限型代数，因此这里使用几何正则性的 Noether 代数版本：设 \(F\) 为域，\(B\) 为 Noether \(F\) 代数；如果对每个有限域扩张 \(F'/F\)，环 \(B\otimes_FF'\) 都正则，就称 \(B\) 在 \(F\) 上几何正则。

\begin{example}\label{b3:exm:formal-fibres-line}
设 \(k\) 为特征零域，\(A=k[x]_{(x)}\)，则 \(\widehat A=k[\![x]\!]\)。\(A\) 的两个素点 \((x)\) 与 \((0)\) 上的形式纤维分别为
\[
k[\![x]\!]\otimes_A k\cong k,
\qquad
k[\![x]\!]\otimes_A k(x)\cong k((x)).
\]
第二式中，对 \(A\) 所有非零元局部化，等价于使 \(x\) 可逆：非零多项式在 \(k[\![x]\!]\) 中是某个 \(x\) 次幂乘单位，反过来 \(x\) 本来就在局部化集合中。于是得到 Laurent 级数域。闭纤维在其基域 \(\kappa((x))=k\) 上几何正则，因为对每个有限扩张 \(L/k\)，基变换后的环就是域 \(L\)。泛形式纤维的基域则是 \(\kappa((0))=k(x)\)。对任意有限扩张 \(L/k(x)\)，特征零保证该扩张可分，所以
\[
k((x))\otimes_{k(x)}L
\]
是有限约化 Artin \(k((x))\) 代数，因而为有限个域之积，仍是正则环。故泛形式纤维在 \(k(x)\) 上也几何正则。
\end{example}

对于交换 Noether 环之间的同态 \(A\to B\)，如果它平坦，且对每个 \(\mathfrak p\in\operatorname{Spec}A\)，纤维 \(B\otimes_A\kappa(\mathfrak p)\) 在 \(\kappa(\mathfrak p)\) 上几何正则，就称该同态为\term[zhengze-huantongtai]{正则环同态}[regular homomorphism]。

\begin{definition}\label{b3:def:C-g-ring}
设 \(A\) 为交换 Noether 环。若对每个 \(\mathfrak p\in\operatorname{Spec}A\)，局部环的极大理想进完备化映射
\[
A_{\mathfrak p}\longrightarrow\widehat{A_{\mathfrak p}}
\]
都是正则环同态，则称 \(A\) 为\term[G-ring]{G 环}[G-ring]。
\end{definition}

\cref{b3:thm:completion-flat}已证明这些完备化映射平坦。因此 G 环条件等价于：每个 \(A_{\mathfrak p}\) 的全部形式纤维在各自的剩余域上几何正则。需要检查的是所有局部环的完备化映射，而不只是某一个局部环的闭纤维。另一个条件是普适链性质：每个有限型 \(A\) 代数中，从一个固定素理想到另一个固定素理想的饱和链具有相同长度。第三个条件通常称为 J-2，要求每个有限型 \(A\) 代数的正则点集为开集。\cref{b3:def:excellent-ring}统一给出这三个条件及卓越环的定义；此处强调它们控制的是不同问题。

\ProseParagraphBreak
域上的有限型代数、它们的局部化以及完备 Noether 局部环属于卓越环；\cref{b3:thm:excellent-basic-classes}给出陈述及证明概要，分别说明链性质、形式纤维及正则点开放性的验证。例如尖点的局部环虽然有奇点，仍是卓越环。因此卓越性并不表示“没有奇点”；它给出研究奇点时所需的一组良好比较性质。在具体问题中，形式纤维、正则点的开放性和链长度条件各有不同的作用。

% !TeX root = ../../Book3.tex
% 纲要标题：### C.2 Frobenius 方法与正特征奇点
% 纲要位置：工作记录/计划纲要.md:2873

\section{Frobenius 方法与正特征奇点}
\label{b3:sec:C02}

% 纲要内容（仅作写作计划，尚非教材正文）：
% * Frobenius 模结构、完美域多项式环的自由基及节点的不平坦算例
% * 迭代分裂映射证明所有理想Frobenius闭；引用Kunz定理说明正则性与Frobenius平坦性的联系
% <!-- 短节：不设小节 -->


% 正文以具体模型说明专题；长证明给出概要并指明技术细节的出处。

设 \(p\) 为素数，\(R\) 为特征 \(p\) 的交换环。二项式系数的消失使 \(F:R\to R,r\mapsto r^p\) 成为环同态，称为\term{Frobenius 同态}[Frobenius homomorphism]。它把所有元素的幂次同时放大，却保留加法。这让正特征中的奇点研究可以把“取幂”转成模论问题。

\begin{definition}\label{b3:def:frobenius-pushforward}
设 \(p\) 为素数，\(R\) 为特征 \(p\) 的交换环，\(e\ge1\) 为整数，\(q=p^e\)。记 \(F_*^eR\) 为加法群仍等于 \(R\)、标量作用改为
\[
a\cdot F_*^e r=F_*^e(a^q r)
\]
的 \(R\) 模，并记 \(F_*R=F_*^1R\)。若 \(F_*R\) 为有限生成的 \(R\) 模，则称 \(R\) 为 \term[finiteF]{F-有限环}[F-finite ring]。若 \(R\) 线性映射 \(R\to F_*R,a\mapsto F_*a^p\) 有 \(R\) 线性左逆，则称 \(R\) 为\term{Frobenius 分裂}[Frobenius split]环。
\end{definition}

标量作用的改变是必要的：普通映射 \(r\mapsto r^p\) 一般不是通常 \(R\) 模之间的线性映射。现在它成为线性的，才可以问平坦、有限或分裂。

\begin{proposition}\label{b3:prop:polynomial-frobenius-free}
设 \(k\) 为特征 \(p>0\) 的完美域，\(R=k[x_1,\ldots,x_n]\)，\(e\ge1\) 为整数，\(q=p^e\)。则 \(F_*^eR\) 为秩 \(q^n\) 的自由 \(R\) 模，一组基为
\[
F_*^e(x_1^{a_1}\cdots x_n^{a_n}),\qquad0\le a_i<q.
\]
特别地，\(R\) 的 Frobenius 同态平坦，且 \(R\) 为 F-有限并 Frobenius 分裂。
\end{proposition}
\begin{proof}
每个指数向量唯一写为 \(qb+a\)，其中各 \(a_i\) 在给定范围内；每个系数在完美域中又有唯一的 \(q\) 次根。因此每个多项式唯一写成 \(\sum_a f_a(x)^q\cdot x^a\)。这正是关于所列基的唯一线性表示，故该模自由。对于 \(e=1\)，把 \(F_*1\) 的系数取出、其余基向量送到零，得到线性映射 \(F_*R\to R\)，它在 \(F_*f^p\) 上的值为 \(f\)，即所需左逆。
\end{proof}

Frobenius 的平坦性实际能够刻画任意正特征 Noether 环的正则性；上述多项式环的自由性是一个可直接计算的例子。这一判据称为\term[Kunz-dingli]{Kunz 定理}[Kunz's theorem]。\footnote{昆茨（Ernst Kunz，1933-2021），德国数学家，研究交换代数与代数几何，引入 Hilbert--Kunz 重数。}原始出处见\cite{Kunz1969RegularLocalRings}。

\begin{claim}[Kunz]\label{b3:claim:C-kunz}
设 \(p\) 为素数，\(R\) 为特征 \(p\) 的交换 Noether 环。则 \(R\) 是正则环，当且仅当 Frobenius 同态
\[
F:R\longrightarrow R,\qquad r\longmapsto r^p
\]
平坦，亦即 \(F_*R\) 是平坦 \(R\) 模。
\end{claim}

这一判据的证明可先归约到局部环，再把 Frobenius 平坦性给出的长度等式与 Hilbert--Samuel 函数的增长次数比较，从而得到极大理想的最少生成元数等于环的维数。

\begin{example}[分裂不等于正则]\label{b3:exm:node-frobenius-split}
设 \(k\) 为特征 \(p>0\) 的完美域，\(R=k[x,y]/(xy)\)。每个元素唯一写为一个常数加只含正次 \(x\) 的多项式与只含正次 \(y\) 的多项式。定义 \(\varphi:F_*R\to R\)：指数可被 \(p\) 整除的单项式保留并把指数除以 \(p\)、系数取 \(p\) 次根，其余单项式送到零。按上述唯一表示检查 \(\varphi(a\cdot F_*r)=a\cdot\varphi(F_*r)\)，只需检查标量为系数、\(x\)、\(y\) 的情形；交叉乘积因 \(xy=0\) 同时为零。且 \(\varphi(F_*a^p)=a\)，故该环 Frobenius 分裂。但原点局部环维数一，而 \(\mathfrak m/\mathfrak m^2\) 维数二，所以不正则。
\end{example}

这一例子还可直接检测 Frobenius 不平坦。若 \(N=F_*R\) 平坦，对映射 \(R\xrightarrow{x}R\) 张量，应有
\[
\ker(x:N\to N)=(0:_Rx)N=yN.
\]
而 \(F_*y\) 被 \(x\) 零化，因为 \(x\cdot F_*y=F_*(x^py)=0\)；它却不在 \(yN=F_*(y^pR)\) 中。于是 Frobenius 虽分裂却不平坦。两种性质不能混为一种正则性判断。

\ProseParagraphBreak
对理想 \(I\subseteq R\)，记 \(I^{[q]}=(a^q:a\in I)\)，并定义\term{Frobenius 闭包}[Frobenius closure]
\[
I^F=\{r\in R:r^{p^e}\in I^{[p^e]}\text{ 对某个 }e\ge0\}.
\]
它是包含 \(I\) 的理想：对两个见证取足够大的共同 \(e\)，用 Frobenius 保持加法与乘法验证。分裂映射则能够把这些取幂后的理想关系还原到原环。

\begin{proposition}\label{b3:prop:C-split-frobenius-closure}
设 \(p\) 为素数，\(R\) 为特征 \(p\) 的交换环。如果 \(R\) 为 Frobenius 分裂环，则每个理想 \(I\subseteq R\) 都满足 \(I^F=I\)。
\end{proposition}
\begin{proof}
取 \(R\) 线性分裂映射 \(\varphi:F_*R\to R\)，即
\[
\varphi(F_*r^p)=r\qquad(r\in R).
\]
置 \(\varphi_1=\varphi\)，并归纳定义
\[
\varphi_{e+1}=\varphi\circ F_*\varphi_e:
F_*^{e+1}R\longrightarrow R.
\]
这里 \(F_*\varphi_e:F_*^{e+1}R\to F_*R\) 是同一加法映射在两端限制标量后所得的映射，即
\[
F_*^{e+1}t\longmapsto
F_*\bigl(\varphi_e(F_*^e t)\bigr).
\]
它仍为 \(R\) 线性映射，因为两端的标量作用都多复合了一次 Frobenius。若 \(\varphi_e(F_*^e r^{p^e})=r\)，则
\[
\varphi_{e+1}(F_*^{e+1}r^{p^{e+1}})
=\varphi\bigl(F_*\varphi_e(F_*^e(r^p)^{p^e})\bigr)
=\varphi(F_*r^p)=r.
\]
故每个 \(\varphi_e\) 都分裂第 \(e\) 次 Frobenius。

若 \(r\in I^F\)，则当见证为 \(e=0\) 时已有 \(r\in I\)；否则取 \(q=p^e\)，将 \(r^q\in I^{[q]}\) 写成有限和
\[
r^q=\sum_{i=1}^s a_i^q b_i,
\qquad a_i\in I,\quad b_i\in R.
\]
由 \(\varphi_e\) 的 \(R\) 线性性得到
\[
r=\varphi_e(F_*^e r^q)
=\sum_{i=1}^s a_i\,\varphi_e(F_*^e b_i)\in I.
\]
因此 \(I^F\subseteq I\)，与反向包含合起来即得结论。
\end{proof}

\ProseParagraphBreak
更精细的正特征奇点理论还研究所有模上的纯性、局部上同调上的 Frobenius 作用及紧闭包。它们分别施加不同的有限性与非零条件。

% !TeX root = ../../Book3.tex
% 纲要标题：### C.3 余切复形、切复形与形变问题
% 纲要位置：工作记录/计划纲要.md:2880

\section{余切复形、切复形与形变问题}
\label{b3:sec:C03}

% 纲要内容（仅作写作计划，尚非教材正文）：
% * 区分展示给出的朴素两项复形与完整余切复形；特征零完全交的标准DG模型采用已核对的引用结果
% * 用Koszul模型解释方程及其关系，导出对偶得到切复形；超曲面命题只直接计算给定展示的方程扰动
% <!-- 短节：不设小节 -->


% 正文以具体模型说明专题；长证明给出概要并指明技术细节的出处。

Kähler 微分通过 \(\mathrm{d}f\) 记录方程的一阶变化，但它是一个商模，已把产生关系的映射压缩掉了。形变问题往往还需要知道关系怎样变化，以及关系之间是否有进一步关系。\chapref{b3:ch:23}的平方零提升计算已经出现了这种需求。

\ProseParagraphBreak
例如设 \(k\) 为特征零域，\(P=k[x]\)、\(I=(x^3)\)、\(B=P/I\)。\(I/I^2\) 是以 \(x^3\) 的类为基的自由 \(B\) 模，\cref{b3:thm:conormal-sequence}的商环微分正合列中，关系映射因而为
\[
B\longrightarrow B\,\mathrm{d}x,
\qquad a\longmapsto 3ax^2\,\mathrm{d}x.
\]
普通微分模只保存这个映射的余核。例如，关系模中的 \(x\cdot\overline{x^3}\) 并未为零，它的微分却为零。若只保存微分模，就无从恢复这部分关系信息。保留关系模及其映射，将它们放在相邻次数上，便得到下面的复形。

\ProseParagraphBreak
设 \(k\) 为域，\(\alpha:P=k[x_1,\ldots,x_n]\twoheadrightarrow B\) 是交换 \(k\) 代数的一个展示，核为 \(I\)。无论 \(I\) 是否由正则序列生成，\cref{b3:thm:conormal-sequence}的商环微分正合列都给出两项复形
\[
NL(\alpha)=\bigl[I/I^2\xrightarrow{\mathrm{d}}\Omega_{P/k}\otimes_PB\bigr],
\]
次数分别为 \(-1,0\)，它称为该展示的\term{朴素余切复形}[naive cotangent complex]。其零次上同调为 \(\Omega_{B/k}\)，负一次上同调保留关系映射的核。真正的\term{余切复形}[cotangent complex] \(L_{B/k}\) 则由代数消解构造，继续记录关系之间的高阶关系；一般情形下，它不能由上述两项代替。

\ProseParagraphBreak
若 \(k\) 为特征零域，\(P=k[x_1,\ldots,x_n]\)、\(I=(f_1,\ldots,f_c)\)，且 \(f_1,\ldots,f_c\) 是 \(P\) 中的正则序列，则对 \(B=P/I\)，完全交余切复形的标准计算给出
\[
L_{B/k}\simeq\bigl[I/I^2\xrightarrow{\mathrm{d}}\Omega_{P/k}\otimes_PB\bigr]
\cong\bigl[B^c\xrightarrow{J}B^n\bigr],
\]
其中 \(J(e_j)=\sum_i(\partial f_j/\partial x_i)\,\mathrm{d}x_i\)，次数仍为 \(-1,0\)。这里使用的是特征零下准自由微分分次代数计算余切复形的理论\footnote{所引文献的全篇约定基域代数闭，而 Definition 2.1 的代数准备先只要求特征零。上述模型与标量扩张相容；由 Section 2.3 的基变换公式 (2.2)，比较映射扩至代数闭包后为拟同构，再由忠实平坦性知它在原基域上也是拟同构。}，见\cite[Section 2.3, Proposition 2.12]{BravBussiJoyce2019Darboux}。下述 Koszul 构造说明正则序列如何使这一计算只剩两项。

\begin{proofsketch}
给 \(P\) 添加次数 \(-1\) 的反交换生成元 \(\xi_1,\ldots,\xi_c\)，令
\[
Q=P\otimes_k\bigwedge(\xi_1,\ldots,\xi_c),
\qquad \mathrm{d}_Q\xi_j=f_j,\qquad \mathrm{d}_Q x_i=0.
\]
正则序列的 Koszul 正合性，见\cref{b3:thm:koszul-regular-resolution}，说明增广映射 \(Q\to B\) 是拟同构。记 \(\delta:Q\to\Omega_{Q/k}\) 为泛导子，以区别于复形微分 \(\mathrm{d}_Q\)。微分模 \(\Omega_{Q/k}\) 在分次意义下自由，由次数零的 \(\delta x_i\) 和次数 \(-1\) 的 \(\delta\xi_j\) 生成；其复形微分将 \(\delta\xi_j\) 送到
\[
\delta f_j=\sum_i\frac{\partial f_j}{\partial x_i}\,\delta x_i.
\]
应用上述余切复形计算结果，\(L_{B/k}\) 由 \(B\otimes_Q\Omega_{Q/k}\) 表示，因而得到 \([B^c\xrightarrow{J}B^n]\)。正则序列还使 \(I/I^2\) 以各 \(f_j\) 的类为基，给出所写的关系模形式。若关系不构成正则序列，Koszul 复形可能有非零负次上同调；一旦出现这种情形，\(Q\to B\) 就不再是所需的消解，须继续添加生成元处理高阶关系。
\end{proofsketch}

低阶余切理论、完全交判据及无穷小形变的论述，见\cite[Sections 2.1, 3.2, 4.3]{LichtenbaumSchlessinger1967CotangentComplex}。

\ProseParagraphBreak
回到三重加厚点，设 \(k\) 为特征零域，并记 \(B_0=k[x]/(x^3)\)，模型就是 \([B_0\xrightarrow{3x^2}B_0]\)。若 \(a=a_0+a_1x+a_2x^2\)，则 \(3x^2a=3a_0x^2\)，故
\[
H^{-1}(L_{B_0/k})=(x)=kx\oplus kx^2,
\qquad H^0(L_{B_0/k})=B_0/(x^2).
\]
两个模都只有二维，却位于不同次数，承担不同作用：零次是普通微分，负一次保留关系映射的核。这也说明，引入余切复形后不能只检查它的零次上同调。

\ProseParagraphBreak
对前面的一般完全交 \(B=P/I\)，余切模型的两项都是有限自由模，所以逐项取对偶已计算其导出对偶。所得
\[
T_{B/k}=\mathbf R\operatorname{Hom}_B(L_{B/k},B)
\simeq\bigl[B^n\xrightarrow{J^{\mathsf T}}B^c\bigr]
\]
位于次数 \(0,1\)，称为\term{切复形}[tangent complex]的模型。按内部 Hom 的约定，其微分原为 \(-J^{\mathsf T}\)；在一次乘以 \(-1\)，便得到右侧所写的同构模型。零次核就是保持全部关系的导子 \(\operatorname{Der}_k(B,B)\)，一次余核则度量不能由坐标变化消掉的方程扰动。余切复形与切复形因此是对偶关系，不是同一记号的两个名称。一般非自由情形须先作消解，再用导出 Hom，不能随意逐项对偶。

\begin{definition}\label{b3:def:first-order-deformation}
设 \(k\) 为域，\(B\) 为交换 \(k\) 代数，\(D=k[\varepsilon]/(\varepsilon^2)\)。若平坦的交换 \(D\) 代数 \(B_\varepsilon\) 连同指定同构 \(B_\varepsilon/\varepsilon B_\varepsilon\cong B\) 一起给出，则称这组数据为 \(B\) 的\term{一阶形变}[first-order deformation]。形变之间的同构须与该指定同构相容。
\end{definition}

平坦性排除了任意添加幂零关系。例如 \(D[x]/(\varepsilon x)\) 在模 \(\varepsilon\) 后也是 \(k[x]\)，却把 \(x\) 在无穷小方向额外零化，因而不是所需的平坦变化。一般定义见\cite[Section 4.3, Definition 4.3.2]{LichtenbaumSchlessinger1967CotangentComplex}。固定一个超曲面方程后，可以直接计算方程的扰动，以及坐标变化能消去哪些扰动。

\begin{proposition}[超曲面的方程扰动]\label{b3:prop:hypersurface-first-order-perturbation}
设 \(k\) 为特征零域，\(P=k[x_1,\ldots,x_n]\)，\(f\in P\) 为非零非单位，\(B=P/(f)\)，\(D=k[\varepsilon]/(\varepsilon^2)\)。对任意 \(g\in P\)，代数
\[
B_g=D[x_1,\ldots,x_n]/(f+\varepsilon g)
\]
连同模 \(\varepsilon\) 后与 \(B\) 的自然同构，构成 \(B\) 的一阶形变。在给定展示 \(B=P/(f)\) 下，把扰动按坐标变换 \(x_i\mapsto x_i+\varepsilon a_i\) 及生成方程乘 \(1+\varepsilon b\) 视为等价，其中 \(a_i,b\in P\)。这些方程扰动的等价类由
\[
P/(f,\partial f/\partial x_1,\ldots,\partial f/\partial x_n)
\]
参数化。令 \(T_{B/k}=[B^n\xrightarrow{(\partial f/\partial x_i)}B]\) 为次数 \(0,1\) 上的切复形模型，则这个参数模也同构于 \(H^1(T_{B/k})\)。
\end{proposition}
\begin{proof}
选取 \(k\) 向量空间补空间 \(V\)，使 \(P=V\oplus fP\)。将 \(h_0+\varepsilon h_1\) 中 \(h_0\) 写成 \(v_0+fq_0\)，再用 \(f=-\varepsilon g\) 代换，并将 \(h_1-gq_0\) 模 \(fP\) 化简，得到唯一形式 \(v_0+\varepsilon v_1\)，其中 \(v_i\in V\)。唯一性如下：若该式等于 \((f+\varepsilon g)(q_0+\varepsilon q_1)\)，比较常数项得 \(v_0=fq_0\)，所以 \(v_0=q_0=0\)；再比较 \(\varepsilon\) 项得 \(v_1=fq_1\)，故 \(v_1=0\)。这里用到 \(P\) 是整环且 \(f\ne0\)。因此 \(B_g\) 作为 \(D\) 模自由，同构于 \(D\otimes_kV\)，故平坦。

平方零 Taylor 公式给出
\[
f(x+\varepsilon a)+\varepsilon\cdot g(x+\varepsilon a)
=f+\varepsilon\left(g+\sum_i a_i\dfrac{\partial f}{\partial x_i}\right).
\]
再乘 \(1+\varepsilon b\) 会给扰动加上 \(bf\)。这些变换都有把 \(\varepsilon\) 项取相反数得到的逆，且其复合仍只对 \(g\) 添加上述理想中的元素。因此等价关系恰是模该理想同余。最后在 \(B\) 中取 Jacobian 行向量的余核，得到同一个商。
\end{proof}

一般形变分类定理进一步说明，以平凡形变 \(D\otimes_kB\) 为零元后，\(B\) 的全部抽象一阶形变同构类由 \(T^1(B/k,B)\) 参数化；见\cite[Section 3.1, Lemma 3.1.2; Section 4.3, Theorem 4.3.3]{LichtenbaumSchlessinger1967CotangentComplex}。这里 \(T^1(B/k,B)\) 表示一次余切上同调，完全交模型将它识别为 \(H^1(T_{B/k})\)。因此，上述商模既可以从方程与坐标变化直接算出，也在这一分类定理下描述抽象一阶形变。

\begin{example}\label{b3:exm:cusp-tangent-complex}
对尖点 \(B=k[x,y]/(y^2-x^3)\)，\(\operatorname{char}k=0\)，上式给出
\[
H^1(T_{B/k})=B/(x^2,y)\cong k[x]/(x^2).
\]
所以任意方程扰动在上述等价下可化为 \(y^2-x^3+\varepsilon(a+bx)\)，且两个参数不能被这种坐标变化消去。若只检查光滑点的切空间，便看不到这一整族方程的变化。
\end{example}

从 \(\varepsilon^2=0\) 延长到更高次幂零参数时，还须检查新的相容性；一般形变的障碍由更高的上同调群表达。上述计算处理的是仿射超曲面的方程扰动。更一般的形变还要考虑关系之间的关系，全局问题则另外需要局部形变的相容黏合。

% !TeX root = ../../Book3.tex
% 纲要标题：### C.4 局部交数、乘数理想与对偶复形
% 纲要位置：工作记录/计划纲要.md:2887

\section{局部交数、乘数理想与对偶复形}
\label{b3:sec:C04}

% 纲要内容（仅作写作计划，尚非教材正文）：
% * 平面曲线有限长度交数与Tor算例；区分不满足有限长度的自交
% * 对数消解存在性采用主理想化定理；乘数理想的解析—代数比较与选择独立性采用准确引用
% * 正规交叉模型中展开消失阶与局部可积性的对应，不以一行完备化说明代替比较定理
% * 复用第22章局部对偶，区分CM的单一次数与非CM商环的多次数Ext
% #### C.4.1 局部交数与 Tor
% #### C.4.2 对数消解与乘数理想
% #### C.4.3 非 Cohen--Macaulay 环的多次数对偶


% 正文以具体模型说明专题；长证明给出概要并指明技术细节的出处。

相交的厚度、函数沿除子的消失阶与自由消解的对偶，分别产生长度、理想和上同调模。它们依赖的对象也不同：交数研究一对曲线，乘数理想研究光滑簇上的理想及其正有理参数，而局部对偶比较局部上同调与 Ext。下面从具体计算说明这些差别。

\subsection{局部交数与 Tor}
\label{b3:subsec:C04:01}

两条平面曲线在一点相遇时，仅记录这个点不够：横截与相切应有不同的交数。局部商环的长度提供最直接的模型。\cref{b3:def:plane-local-intersection-number}给出复平面曲线的定义；这里同样设 \(k\) 为域，\(S=k[\![x,y]\!]\)，曲线由 \(f,g\) 定义，并假定 \(S/(f,g)\) 为有限长度模；在这个平面模型中，把
\[
i_0(f,g)=\ell_S(S/(f,g))
\]
称为原点处的\term{局部交数}[local intersection multiplicity]。

\begin{example}\label{b3:exm:plane-contact-multiplicity}
设 \(k\) 为域，\(S=k[\![x,y]\!]\)。对 \(f=y\)、\(g=y-x^m\)，其中 \(m\ge1\)，有
\[
S/(f,g)\cong k[\![x]\!]/(x^m),\qquad i_0(f,g)=m.
\]
模的滤过由 \(1,x,\ldots,x^{m-1}\) 给出 \(m\) 个长度一因子。另一方面，用 \(0\to S\xrightarrow{y}S\to S/(y)\to0\) 与 \(S/(y-x^m)\) 张量，乘 \(y\) 化为 \(k[\![x]\!]\) 上乘 \(x^m\)，是单射。因此
\[
\operatorname{Tor}_1^S(S/(y),S/(y-x^m))=0,
\]
而零次张量积的长度为 \(m\)。这把交数与本卷的自由消解计算联系起来。
\end{example}

若改为两次同一条曲线 \(y=0\)，零次张量积与一次 Tor 都同构于 \(k[\![x]\!]\)，不再具有有限长度。此时不能写下两个无限长度的形式差来定义交数。一般交理论使用循环、等价关系及更系统的 Tor 交替和，需要额外的有限性和维数假设；本节只证明以上有限长度模型。

\subsection{对数消解与乘数理想}
\label{b3:subsec:C04:02}

另一类信息来自函数沿除子的消失阶。为说明乘数理想，先固定光滑不可约复代数簇 \(X\)、非零凝聚理想层 \(\mathfrak a\) 和有理数 \(c>0\)。选取一个\term{对数消解}[log resolution] \(\mu:Y\to X\)：它是光滑簇之间的固有双有理态射，使 \(\mathfrak a\mathcal O_Y=\mathcal O_Y(-F)\)，且有效除子 \(F\) 与例外除子的支集具有简单正规交叉，即局部由若干坐标超平面的并给出。\(\mathcal O_Y(D)\) 表示局部允许的有理函数层，沿每个素除子的极点阶不超过 \(D\) 的相应系数。相对典范除子 \(K_{Y/X}\) 由 Jacobian 行列式记录顶次微分形式的变化。于是定义\term{乘数理想}[multiplier ideal]
\[
\mathcal J(\mathfrak a^c)
=\mu_*\mathcal O_Y\bigl(K_{Y/X}-\lfloor cF\rfloor\bigr),
\]
其中下取整逐除子系数进行。它首先是在 \(Y\) 上规定函数应消失多少阶，再把这些条件收回 \(X\)。定义见\cite[Lecture 1, Definition 1.3]{Lazarsfeld2010ShortCourse}；所用对数消解的存在性来自特征零主理想化定理，其选择独立性则由双有理变换公式保证。

\begin{theorem}\label{b3:thm:log-resolution-multiplier-independence}
设 \(X\) 为光滑不可约复代数簇，\(\mathfrak a\subseteq\mathcal O_X\) 为非零凝聚理想层，\(c>0\) 为有理数。\(\mathfrak a\) 有对数消解，且
\(\mu_*\mathcal O_Y(K_{Y/X}-\lfloor cF\rfloor)\) 不依赖所选对数消解 \(\mu\)。
\end{theorem}
\begin{proofsketch}
存在性采用特征零的主理想化定理：经过有限次沿光滑中心的吹起，理想的总变换成为简单正规交叉除子的可逆理想，且与例外除子的支集同时具有正规交叉。该结果及理想层版本见\cite[Theorem 1.10, Remark 1.18]{BierstoneMilman1997CanonicalResolution}。构造根据局部不变量选择与已有例外除子正规交叉的光滑中心，吹起后用变换后的不变量继续选择中心，直至理想成为所需的单项式形状。中心选择的相容性与有限终止所用的不变量性质，是\cite[Theorem 1.14, Remark 1.18]{BierstoneMilman1997CanonicalResolution}所建立的深层结论。

先说明所定义的推前层确实为 \(\mathcal O_X\) 的子层。设 \(u\) 为它的一个有理截面。在 \(X\) 任一素除子的泛点附近，固有双有理态射 \(\mu\) 为同构，相对典范除子的系数为零；相应截面条件给出
\(\operatorname{ord}_D(u)\ge\lfloor c\operatorname{ord}_D(\mathfrak a)\rfloor\ge0\)。所以 \(u\) 在 \(X\) 的所有余维一处无极点。光滑 \(X\) 正规，其正则函数环为这些高度一赋值环的交，故 \(u\) 是正则函数。

独立性是\cite[Lecture 1, Proposition 1.6]{Lazarsfeld2010ShortCourse}的结论。用解析方法解释时，要同时采用该文 Proposition 1.8 的解析—代数比较定理：按局部可积性定义的解析乘数理想，是上述代数乘数理想的解析化；其除子形式的双有理变换公式见 Lemma 1.9。下面的坐标计算说明可积条件怎样给出定义中的除子系数。

设 \(f_1,\ldots,f_r\) 为 \(\mathfrak a\) 的局部生成元。对于一个正则函数 \(u\)，考虑
\[
\frac{|u|^2}{(|f_1|^2+\cdots+|f_r|^2)^c}
\]
是否在各点附近可积。更换有限生成集只把分母乘上局部上下有正界的因子，故该条件不变。在任一对数消解的坐标图中，\(f_j\circ\mu\) 的共同因子是 \(z_1^{a_1}\cdots z_l^{a_l}\)，剩余因子共同生成单位理想；Jacobian 的消失阶 \(b_i\) 正是 \(K_{Y/X}\) 的系数。换元后，被积式相当于
\( |u\circ\mu|^2\cdot\prod_i|z_i|^{2(b_i-ca_i)}\)。

一变量积分 \(\int_0^\varepsilon r^{2t+1}\,\mathrm{d}r\) 有限当且仅当 \(t>-1\)。逐变量应用，得到可积性恰好要求
\[
\operatorname{ord}_{z_i=0}(u\circ\mu)
\ge\lfloor ca_i\rfloor-b_i\qquad\text{对所有 }i.
\]
这正是 \(u\) 属于所写推前理想的条件。可积性是在 \(X\) 的解析邻域上定义的，不依赖消解选择；借助上述解析—代数比较与双有理变换定理，便得到代数乘数理想的选择独立性。
\end{proofsketch}

\begin{example}\label{b3:exm:snc-multiplier-ideal}
设 \(c>0\) 为有理数，\(1\le r\le n\)。在 \(X=\mathbb A_{\mathbb C}^n\) 上，设 \(\mathfrak a=(x_1^{a_1}\cdots x_r^{a_r})\)，\(a_i\ge1\)。坐标超平面已为简单正规交叉，恒等映射就是对数消解，且相对典范除子为零，故
\[
\mathcal J(\mathfrak a^c)
=\bigl(x_1^{\lfloor ca_1\rfloor}\cdots x_r^{\lfloor ca_r\rfloor}\bigr).
\]
例如 \(\mathfrak a=(x^2y^3)\)、\(c=1/2\) 时得到 \((xy)\)。这个公式针对一个单项式生成的主理想，不能直接把一个多生成元单项式理想中的各指数分别取整；后者需要考虑整体指数集合。
\end{example}

要让相对典范除子实际参与计算，可取 \(X=\mathbb A_{\mathbb C}^2\)、\(\mathfrak a=(x,y)\)，并吹起原点。第一张图为 \(x=u,y=uv\)，于是
\[
(x,y)\mathcal O_Y=(u),\qquad
\mathrm{d}x\wedge\mathrm{d}y=\mathrm{d}u\wedge(v\,\mathrm{d}u+u\,\mathrm{d}v)=u\,\mathrm{d}u\wedge\mathrm{d}v.
\]
第二张图 \(x=rw,y=w\) 给出同样的消失阶。例外除子 \(E\) 在第一张图由 \(u=0\) 定义，故理想的总变换为 \(\mathcal O_Y(-E)\)，而 Jacobian 说明 \(K_{Y/X}=E\)。这是一个对数消解，因而
\[
\mathcal J((x,y)^c)
=\mu_*\mathcal O_Y\bigl((1-\lfloor c\rfloor)E\bigr)
=(x,y)^{\max\{\lfloor c\rfloor-1,0\}}.
\]
最后一式可按消失阶核验：一个底空间正则函数 \(g\) 拉回后沿 \(E\) 的阶数，就是它在原点的最低总次数，因此要求 \(\operatorname{ord}_E(g\circ\mu)\ge\lfloor c\rfloor-1\)。当右侧非正时，对正则函数没有新增条件；推前截面在底空间正则已由上面的余维一论证保证。例如 \(c=2\) 时理想为 \((x,y)\)，而非 \((x,y)^2\)。这个一次的差正来自顶次微分形式中的 Jacobian 因子：理想阶数与相对典范阶数必须相减。

\subsection{非 Cohen--Macaulay 环的多次数对偶}
\label{b3:subsec:C04:03}

\chapref{b3:ch:22}的局部对偶给出另一种信息。给定 \(n\) 维完备正则局部环 \(S\) 与非零商 \(R=S/I\)，\cref{b3:thm:regular-local-duality}已经给出
\[
D_S\bigl(H_{\mathfrak m}^i(R)\bigr)
\cong\operatorname{Ext}_S^{n-i}(R,S).
\]
若 \(R\) 为维数 \(d\) 的 CM 环，右侧仅在 \(n-d\) 次非零，典范模便足以承载信息。非 CM 时可能有多个非零次数；\cref{b3:def:regular-cover-dual-complex}用整个对偶复形把它们保存下来。低次局部上同调需要相应次数的 Ext 与之对偶，整个复形还保留这些次数之间的联系。这里的模型以给定正则局部覆盖为基础；一般环上对偶复形的存在性属于更进一步的对偶理论。



\begin{chaptersummary}
奇点研究中的附加结构各有明确任务。形式纤维比较完备化前后的局部信息；Frobenius 把取幂变成线性模结构；切复形记录坐标变化能否消掉方程扰动。交数与乘数理想分别从有限长度和消失阶数提取数值或理想，而非 CM 环的局部对偶要求保存多个次数。

\ProseParagraphBreak
这些构造依赖不同的条件：Frobenius 的自由性计算使用特征与完美性，两项余切模型限于完全交，交数要求有限长度，乘数理想使用光滑复簇及对数消解。主理想化把消失阶数化为正规交叉除子上的系数；代数与解析乘数理想的比较定理把它识别为内在的可积条件，从而解释选择独立性。一般形变与非 Cohen--Macaulay 情形的对偶都需要保留相应的复形信息。
\end{chaptersummary}

\BookExercises

\begin{exercise}\label{b3:ex:C-complete-formal-fibre}
设 \((A,\mathfrak m)\) 为完备 Noether 局部环。计算映射 \(A\to\widehat A\) 在任意素点 \(\mathfrak p\) 的形式纤维。
\end{exercise}
\begin{solution}
完备性给出 \(\widehat A=A\)，故形式纤维为 \(A\otimes_A\kappa(\mathfrak p)=\kappa(\mathfrak p)\)。这只计算本局部环的完备化映射；要按 G 环定义处理全部 \(A_{\mathfrak q}\) 的完备化，还需一般定理，不能由这一行计算直接取代。
\end{solution}

\begin{exercise}\label{b3:ex:C-plane-formal-fibre}
设 \(k\) 为特征零域，\(F=k(\!(y)\!)\)。取 \(k(y)\) 的二次扩张
\[
L_1=k(y)(\sqrt y),\qquad L_2=k(y)(\sqrt{1+y}).
\]
分别计算 \(F\otimes_{k(y)}L_1\)、\(F\otimes_{k(y)}L_2\)。说明即使两者都正则，基变换后的环也不一定是域；联系二维非闭点形式纤维的计算解释这一现象。
\end{exercise}
\begin{solution}
\(y\) 与 \(1+y\) 在有理函数域中分别有奇数阶零点，所以均非平方，两扩张确为二次可分扩张。在 \(F\) 中，平方的 \(y\)-赋值为偶数，故 \(T^2-y\) 仍不可约，第一张量积为二次域 \(F(\sqrt y)\)。另一方面，\(T^2-(1+y)\) 的模 \(y\) 单根 \(\pm1\) 各自提升为 \(\pm u\in k[\![y]\!]\)，且 \(2u\) 为单位。中国剩余定理给出第二张量积 \(F\times F\)。二者都是有限个域的乘积，所以正则；正则性没有要求只有一个域因子。\cref{b3:lem:C-plane-formal-fibre}中的纤维正是以 \(k(y)\) 为底域的 \(F\)，检查几何正则性时须允许这种分裂。
\end{solution}

\begin{exercise}\label{b3:ex:C-frobenius-basis}
在 \(R=\mathbb F_3[x,y]\) 中给出 \(F_*R\) 的一组自由基，并分解 \(F_*(x^7y^5+x^3+2)\)。
\end{exercise}
\begin{solution}
基为 \(F_*(x^ay^b)\)，\(0\le a,b<3\)，共九个。由于 \(x^7y^5=(x^2y)^3xy^2\)、\(x^3=x^3\cdot1\)、\(2=2^3\)，有
\[
F_*(x^7y^5+x^3+2)=x^2y\cdot F_*(xy^2)+(x+2)\cdot F_*1.
\]
模结构中的标量先取三次幂，所以系数是 \(x^2y\) 与 \(x+2\)，不能写成原单项式的普通系数。
\end{solution}

\begin{exercise}\label{b3:ex:C-frobenius-closure-cusp}
设 \(k\) 为特征 \(p>0\) 的域，\(R=k[t^3,t^4,t^5]\)、\(I=(t^3)\)。证明 \(t^4\notin I\)，并求使 \((t^4)^{p^e}\in I^{[p^e]}\) 的最小正整数 \(e\)。据此判断 \(R\) 能否 Frobenius 分裂。
\end{exercise}
\begin{solution}
环 \(R\) 中的单项式指数为 \(0\) 及所有不小于 \(3\) 的整数。若 \(t^4=t^3a\)，在 \(k[t]\) 中消去得 \(a=t\notin R\)。令 \(q=p^e\)，则 \(I^{[q]}=(t^{3q})\)，所以所求成员关系等价于 \(t^q\in R\)，也就是 \(q\ge3\)。因此 \(p=2\) 时最小 \(e=2\)，\(p\ge3\) 时最小 \(e=1\)。这个非闭理想与\cref{b3:prop:C-split-frobenius-closure}相矛盾，故 \(R\) 不能 Frobenius 分裂。
\end{solution}

\begin{exercise}\label{b3:ex:C-dual-numbers-cotangent}
设 \(k\) 为特征零域，\(B=k[x,y]/(x^2,y^2)\)。证明所给两方程为正则序列，写出两项余切复形及切复形，计算 \(H^{-1}(L)\)、\(H^0(L)\)、\(H^0(T)\)、\(H^1(T)\) 的 \(B\) 模结构与 \(k\) 维数。
\end{exercise}
\begin{solution}
\(x^2\) 在多项式整环中为非零因子，而 \(k[x,y]/(x^2)\) 在 \(k[y]\) 上自由，故乘 \(y^2\) 仍单射。两项模型的微分为 \(\operatorname{diag}(2x,2y)\)，余切的次数为 \(-1,0\)，切的次数为 \(0,1\)。因 \(2\) 可逆，\(\operatorname{ann}_B(x)=(x)\)、\(\operatorname{ann}_B(y)=(y)\)，于是
\[
H^{-1}(L)\cong H^0(T)\cong(x)\oplus(y),\qquad
H^0(L)\cong H^1(T)\cong B/(x)\oplus B/(y).
\]
这里 \((x)\) 的基为 \(x,xy\)，\((y)\) 的基为 \(y,xy\)，两个商各为二维；四个上同调模的 \(k\) 维数均为 \(4\)。导出 Hom 可以逐项计算，因为模型的两个项均为有限自由 \(B\) 模。
\end{solution}

\begin{exercise}\label{b3:ex:C-node-deformation}
设 \(k\) 为特征零域，\(f=xy(x-y)\in k[x,y]\)，\(B=k[x,y]/(f)\)。求固定多项式展示下的一阶方程扰动参数空间，给出每个扰动 \(f+\varepsilon g\) 的唯一四参数代表，并比较参数维数与三条分支的数目。
\end{exercise}
\begin{solution}
偏导数为 \(f_x=2xy-y^2\)、\(f_y=x^2-2xy\)。Euler 等式 \(3f=xf_x+yf_y\) 说明 \(f\) 已属于 Jacobian 理想，故参数空间为
\[
k[x,y]/(2xy-y^2,\ x^2-2xy).
\]
在此商中 \(x^2=y^2=2xy\)，又有 \(x^2y=2xy^2=4x^2y\)，所以 \(x^2y=xy^2=0\)，所有三次项均为零。两条独立的二次关系使二次部分恰为由 \(xy\) 张成的一维空间；因此 \(1,x,y,xy\) 是基。每个扰动类唯一表示为 \(a+bx+cy+dxy\)，即由坐标变化及方程生成元重标后可取 \(f+\varepsilon(a+bx+cy+dxy)\)。三个不同线性因子给出三条分支，但扰动空间维数为四，二者是不同信息。
\end{solution}

\begin{exercise}\label{b3:ex:C-intersection-two-thicknesses}
设 \(k\) 为域，\(S=k[\![x,y]\!]\)。求 \(S/(x^a,y^b)\) 的长度，其中 \(a,b\ge1\)。
\end{exercise}
\begin{solution}
商有 \(k\) 基 \(x^iy^j\)，其中 \(0\le i<a,0\le j<b\)，故维数为 \(ab\)。其唯一简单模为剩余域 \(k\)，有限维模的长度等于 \(k\) 维数，因此长度为 \(ab\)。这记录了两条非约化坐标曲线的厚度，不能仅按一个交点计数。
\end{solution}

\begin{exercise}\label{b3:ex:C-multiplier-jumps}
在 \(\mathbb A^2_{\mathbb C}\) 中，令 \(\mathfrak a=(x^2y^3)\)。求 \(c=1/4,1/3,1/2,1\) 时的乘数理想，并确定它第一次不等于全环的正参数。
\end{exercise}
\begin{solution}
使用正规交叉公式，四个结果依次为 \((1),(y),(xy),(x^2y^3)\)。当 \(0<c<1/3\) 时 \(2c,3c\) 都小于一，两个下取整为零；在 \(c=1/3\) 时第二个下取整首次为一。因此第一次变化发生在 \(1/3\)。
\end{solution}

\begin{exercise}\label{b3:ex:C-cm-and-dual-degrees}
设 \(S\) 为 \(n\) 维完备正则局部环，\(R=S/I\ne0\) 为 \(d\) 维 CM 环。由本卷已证定理说明 \(\operatorname{Ext}_S^j(R,S)\) 的非零次数，并解释删去 CM 条件后为何不能照用此结论。
\end{exercise}
\begin{solution}
对有限生成模，深度由局部上同调的首次非零次数刻画，维数以上的局部上同调则消失；这两个结论给出 \(H_{\mathfrak m}^i(R)=0\) 除非 \(i=d\)，且 \(H_{\mathfrak m}^d(R)\ne0\)。正则局部对偶及 Matlis 对偶的忠实性因此给出 \(\operatorname{Ext}_S^j(R,S)\) 仅在 \(j=n-d\) 非零。若非 CM，则深度小于维数，首次非零上同调出现在较低次数，不能由唯一的顶部典范模替代；第22章的非 CM 算例已明确算出两个非零 Ext 次数。
\end{solution}
